\documentclass{article}
\usepackage[T1]{fontenc}
\usepackage{lmodern}
\usepackage{graphicx}
\usepackage{amsmath,amssymb}
\usepackage[a4paper,margin=2cm]{geometry}
\usepackage[absolute,overlay]{textpos}
\setlength{\TPHorizModule}{1mm}
\setlength{\TPVertModule}{1mm}
\usepackage[indonesian]{babel}
\usepackage{hyperref}
\setlength{\emergencystretch}{2em}
% unit: Halaman 51

\begin{document}

% === Halaman 51 (python/native) ===

51

\# Program utama (pemanggilan fungsi enkripsi dan dekripsi AES)

file\_plaintext = input("Ketikkan nama file input pesan yang akan dienkripsi: (beserta path) ")

file\_ciphertext = input("Ketikkan nama file output pesan hasil enkripsi (beserta path): ")

kunci = input("Kunci?: ")

encrypt\_file(file\_plaintext, file\_ciphertext, kunci)

print("Enkripsi selesai. Hasil disimpan di:", file\_ciphertext)

file\_ciphertext = input("Ketikkan nama file input pesan yang akan didekripsi: (beserta path) ")

file\_plaintext = input("Ketikkan nama file output pesan hasil dekripsi (beserta path): ")

kunci = input("Kunci?: ")

decrypt\_file(file\_ciphertext, file\_plaintext, kunci)

print("Dekripsi selesai. Hasil disimpan di:", file\_plaintext)

\end{document}
